%%%%%%%%%%%%%%%%%%%%%%%%%%%%%%%%%%%%%%%%%%%%%%%%%%%%%%%%%%%%%%%%%%%%%%%%%%%%%%%
% *** Kapitelreferenzen ***
%%%%%%%%%%%%%%%%%%%%%%%%%%%%%%%%%%%%%%%%%%%%%%%%%%%%%%%%%%%%%%%%%%%%%%%%%%%%%%%
%
\section{Kapitelreferenzen}\label{sec:kapitelreferenzen}
%
\subsection{Subsection}\label{subsec:subsection}
%
\subsubsection{Subsubsection}\label{subsubsec:subsubsection}
%
\paragraph{Paragraph}\label{paragraph:paragraph}
$\;$ \\ % haesslicher kleiner Trick zum erzwingen des Zeilenumbruchs
\\
Nummerierungstiefe der Kapitel kann in ./formatforlagen/layout.tex bearbeitet werden.\\
Tiefer als Paragraph sollte nicht strukturiert werden.\\\\
%
Referenzierung auf chapter: \kap{chap:formale_elemente}\\\\
Referenzierung auf section: \kap{sec:kapitelreferenzen}\\\\
Referenzierung auf subsection: \kap{subsec:subsection}\\\\
Referenzierung auf subsubsection: \kap{subsubsec:subsubsection}\\\\
Referenzierung auf paragraph: \kap{paragraph:paragraph} (funktioniert nur eindeutig bis zur definierten Nummerierungstiefe)\\\\
%
\newpage
%
